\chapter{設計とトレードオフ}

言語は判断の集合であり、判断には利益だけでなく代償もあります。本章では wick の判断を率直に論じます。同じ実際のゲームの wick 版と Lua 版を並べ、今でも Lua の方が読みやすい点を明言し、v0.3 の省略が恥ずべき欠落ではなく意図的な選択であることを支える、ロードマップの考え方を説明します。

\section{同じゲームを二つの言語で}

Lantern Night がその証拠です。Lua の `games/showcase` と wick の `games/showcase_wick` という、意図的に別々にした二つのプロジェクトとして、両言語で提供しています。同じ実際のゲームで比較し続けられるよう、並べて維持しています。二つをロゼッタストーンのように読み比べてください。三つの抜粋で議論を示します。

\subsection{最高得点の読み込み}

第4章のバグが、本来の場所にあります。Lua 版は、最初のフレームでクラッシュする状態で出荷されました。

\begin{lstlisting}[language=luatwin]
-- Lua: a 0-byte save makes load_save return "" (truthy!), the
-- fallback never fires, and tonumber("") = nil poisons everything.
local best = tonumber(lt.load_save("lanternnight_best") or "0")
\end{lstlisting}

\noindent
wick 版ではそのバグを表現できません。ファイルがない場合と空の場合という、値がない二つの可能性をどちらも処理する必要があるため、下の行だけがコンパイルできます。

\begin{lstlisting}[language=wick]
// wick: the ONLY version the compiler accepts.
let best = num(lt.load_save("lanternnight_wick_best") ?? "") ?? 0
\end{lstlisting}

\noindent
ここでは wick が明確に優れています。安全なコードはバグのある Lua のコードより短く、安全でないコードはコンパイルできません。

\subsection{必要性が示されたレコード}

Wick 0.1 は、ランプや小道具をスカラーの並行リストで表していました。コンパイラを小さく保てる一方、一つの論理オブジェクトが複数のリストに散らばります。KORA のキッチンで、その代償が具体的になりました。五つのリストを同期させる必要があったのです。そのため Wick 0.2 は平坦なレコードを取り入れました。第5章で、検査されるフィールドとリストへの格納を説明します。Lua は今でも、より豊かな入れ子のテーブルとメソッドを扱えます。Wick はその範囲すべてに対応していると主張しません。

Wick 0.3 も同じ規則に従います。プロセッサを組み立てるプロジェクトには、マスク、シフト、レジスタの折り返し、読みやすいバスの値が必要です。Bit Lab はその操作を直接使い、第9章はその成果を説明します。新しい関数は既存の型付きネイティブ関数の仕組みを使い、式の優先順位を変更しません。

\subsection{複数の戻り値}

こちらは小さな不便です。Lua は最も近いランプとその距離を一緒に返します。

\begin{lstlisting}[language=luatwin]
local i, dist2 = nearest_unlit()
\end{lstlisting}

\noindent
wick には複数の戻り値がないので、二つの問い合わせに分け、「ランプがない」は第二の値の欠落ではなく、番兵値の添字にします。

\begin{lstlisting}[language=wick]
let i = nearest_unlit()      // returns the index, or -1
let d = dist2_to(i)          // second query, explicit
\end{lstlisting}

\noindent
記述は長くなりますが、本質的に悪くなったわけではありません。分割は機械的で、番兵値はよく知られた書き方です。複数の戻り値はロードマップにあります。今ないことは不便であって、危険ではありません。

\section{書かなくなるもの}

そうした代償と引き換えに、動的な言語が要求する防御的なコードは実際に減ります。Lantern Night を wick に移植すると、以下はそのままソースから消えました。

\begin{itemize}
\item 保存の読み込みや解析を囲む `if x ~= nil` の連なり。コンパイラが値の欠落の可能性を追跡するので、境界で一度 `??` を書けば、その後は普通の `T` です。
\item `math.` の接頭辞。`sin`、`floor`、`rand` などは名前空間のない組み込み関数です。
\item 整数を末尾の小数なしで表示するためだけの `string.format`。`str(n)` は整数をすっきり表示します。
\item エンジン呼び出しの引数順を過度に心配すること。引数の数と型をコンパイル時に検査し、メッセージに呼び出し名と引数番号が出るので、画面がおかしく見えてからではなく、書いた瞬間に誤りが分かります。
\end{itemize}

\section{変わらないもの}

言語を変えても影響しなかった点を明確にしておきます。wick を採用するコストの範囲が分かるからです。エンジン API は呼び出しごとに同一で、`lt.draw(...)` は両方で同じ行です。`update(dt)` と `draw()` というフレームの契約も同じです。ホットリロードとエラー画面も同じように動きます。そして二つのゲームは同じフレームを描画します。これは比喩ではなく、第6章の決定的なスクリーンショットテストで、wick 版と Lua 版に同じピクセルを要求しています。wick が変えたのは言語で、エンジンや開発の流れではありません。

\section{v0.3 の制約}

wick は意図的に Lua 1.0 のような小ささにしています。以下の省略はすべて、現在の回避策とロードマップ上の位置づけを持つ既知の欠落です。後から発覚する驚きではありません。v0.3 の基準は明確で狭く、実際の lantern ゲームを動かし、Lua の種類のバグをなくし、約二千行を保つことでした。

\begin{center}
\begin{tabular}{@{}p{0.42\linewidth}p{0.50\linewidth}@{}}
\toprule
\textbf{v0.3 にないもの} & \textbf{現在の回避策} \\
\midrule
クロージャ、入れ子・第一級の関数 & トップレベルの \texttt{fn} とモジュールのグローバル変数 \\
複数の戻り値 & 二つの関数、または戻り値の外の状態 \\
入れ子のコンテナ（\texttt{list<list<num>>}） & 平坦なリストの添字計算 \\
モジュール・インポート（\texttt{main.wick} 一つ） & まだアプリ基盤ではなく、ゲームのスクリプト \\
\texttt{for x in list} の反復 & \texttt{for i in 0..len(xs)} \\
文字列の添字・スライス・メソッド & HUD には \texttt{len}、\texttt{+}、\texttt{str()} で対応 \\
戻り値を返す経路の完全な検査 & 型付き関数が末尾まで進むと \texttt{nil}。必ず返すこと \\
グローバル変数・\texttt{and} の連なりの絞り込み & \texttt{??}、または先にローカル変数へコピー \\
指数表記の数値リテラル & 10進表記。16進・2進は対応済み \\
\bottomrule
\end{tabular}
\end{center}

\noindent
二つの点は別に述べる必要があります。これは制約ではなく選択であり、ロードマップでも覆さない判断だからです。

\paragraph{JIT を使わない。} 仕組み上 iOS で許可され、ウォームアップなしに決定的な性能を得られます。これは恒久的な特徴です。

\paragraph{時刻をシードにした乱数を使わない。} 再生と CI のスクリーンショットを可能にするため、決定性を既定にします。実行ごとに変化が必要なら、記録できる数値で自分で `srand()` を設定します。これも恒久的です。

\section{ロードマップの考え方}

次を決める規則は、この言語が従う核心なので、原則として明記します。\emph{wick は本当に必要なものだけを持ち、持つこと自体を目的に加えません。} 実際のゲームのコードで、それがないことによる痛みが示されたときに初めて機能を取り入れます。対になった Lantern Night の二つの版は、常設の証拠の仕組みです。言語を大きくする前に、実際のゲームか焦点を絞った例で、機能の欠落による痛みが現れる必要があります。オプショナル値が本書にあるのは、出荷したゲームのクラッシュが要求したからです。クロージャがないのは、出荷したプログラムがまだその不在の代償を払っていないからです。他の言語と機能をそろえることは、追加の理由にはしません。

全体をまとめる考えは「Lua 1.0 のように小さく」です。Lua はメタテーブル、コルーチン、エコシステムから始まったのではありません。小さく、明快で、組み込みやすい核から始まり、実際の使用が各追加を正当化するにつれて成長しました。wick は意図的に同じ道を進みます。表にある省略はすべて、既存の型付きフロントエンドの背後に追加できます。クロージャと複数の戻り値はコンパイラを考え直す必要はなく、拡張で対応できます。どれも実際のゲームが、規模を増やすだけの必要性を示したときに加え、それ以前には加えません。

だからこそ制約を率直に述べます。欠落を隠す言語は、最悪の時点でユーザーが発見するよう招いていることになります。wick は一覧にし、それぞれに回避策を示し、ロードマップに置きます。プログラムの失敗のしかたは本番で発見するのではなく、事前に見えるべきだという思想が、言語全体の主張だからです。言語は、ユーザーのコードに求めるのと同じ率直さを、ユーザーへ示す責任があります。
